% =====================================================================
% FaaS-MADiG: price-free greedy-diffusion variant of FaaS-MADeA.
% This file is meant to be \input{} (or pasted) into the paper.
% It assumes the paper's notation (N, F, C_i^f, rho_i, omega_i^f,
% beta_{ij}^f, gamma_i^f, p_i^f, ...) is already defined.
% Cross-references are written as plain text "Eq.~(7)", "Alg.~4", etc.;
% convert them to \ref{} once the labels of the host paper are known.
% =====================================================================

\section{FaaS-MADiG: a price-free greedy-diffusion variant}
\label{sec:madig}

\textsc{FaaS-MADiG} (Multi-Agent Distributed Greedy) is a \emph{price-free}
ablation of the \textsc{FaaS-MADeA} auction. It leaves the decentralized
per-iteration loop unchanged: at each control period $t$ every node $i$ still
solves its local problem~(P2) to fix the replicas $r_i^f$, the locally executed
load $x_i^f$, and the residual demand $\omega_i^f(h)$ to offload; it still
advertises on the neighbour-accessible blackboard its residual execution
capacity $C_i^f(h)=\max\{0,\mathrm{Cap}_i^f(h)-x_i^f(h)\}$~(Eq.~(12)) and its
memory slack $\rho_i(h)$~(Eq.~(13)); and it still re-solves the restricted
problem after coordination. \textsc{FaaS-MADiG} replaces \emph{only} the market
layer---bidding, pricing, and the price update of Eqs.~(7)--(8) and
Algorithms~2--3---with a deterministic greedy diffusion that consumes the same
one-hop residual-capacity announcements. Its purpose is to isolate the
contribution of the price signal: \emph{does congestion pricing improve the
allocation over a plain greedy rule that uses exactly the same local
information?}

% =====================================================================
% BEGIN self-contained notation recap.
% Reproduces the relevant rows of the paper's Table 2 (FRALB) and Table 3
% (auction) plus Eqs. (1),(2),(12),(13), so the note is readable on its own.
% REMOVE this whole subsection when inserting into the paper, where the
% notation and these equations are already defined.
% =====================================================================
\subsection{Notation and capacity model}
\label{sec:madig-notation}

Each control period $t$ runs auction iterations $h=0,1,\dots,H_{\max}$ in
parallel. Before coordination, every node $i$ solves a local problem~(P2)
that, for each function $f\in\mathcal{F}_i$, splits the incoming workload
$\lambda_i^f$ into locally processed load $x_i^f$, horizontally offloaded
demand $\omega_i^f$, and Cloud-forwarded load $z_i^f$, and fixes the number of
replicas $r_i^f$. Replicas obey the memory budget and provide capacity
\begin{equation}
  \sum_{f\in\mathcal{F}_i} r_i^f\,\mathrm{RAM}^f \le \mathrm{RAM}_i,
  \qquad
  \mathrm{Cap}_i^f = \frac{r_i^f\,U^f_{\max}}{D_i^f}\ \text{[req/s]} ,
  \tag{1,2}
\end{equation}
from which each node advertises, on a neighbour-accessible blackboard, its
residual execution capacity and its memory slack:
\begin{equation}
  C_i^f(h) = \max\!\big\{0,\ \mathrm{Cap}_i^f(h) - x_i^f(h)\big\},
  \qquad
  \rho_i(h) = \mathrm{RAM}_i - \sum_{f\in\mathcal{F}_i} r_i^f(h)\,\mathrm{RAM}^f \ge 0 .
  \tag{12,13}
\end{equation}
Table~\ref{tab:madig-notation} collects the symbols used below.

\begin{table}[t]
\centering
\caption{Notation used in this section (subset of the paper's Tables~2--3).}
\label{tab:madig-notation}
\small
\begin{tabular}{@{}l l p{0.62\linewidth}@{}}
\toprule
\multicolumn{3}{@{}l}{\textit{Sets}}\\
Nodes / neighbours & $\mathcal{N}$, $N_i$ & computing nodes; one-hop neighbours of $i$.\\
Functions & $\mathcal{F}_i$ & functions deployable on $i$; $\mathcal{F}=\bigcup_i \mathcal{F}_i$.\\
Candidate sellers & $A_i^f(h)$ & neighbours of $i$ with positive score for $f$, $A_i^f(h)\subseteq N_i$.\\
\midrule
\multicolumn{3}{@{}l}{\textit{Parameters and weights}}\\
Input workload & $\lambda_i^f$ & incoming request rate for $f$ at $i$ [req/s].\\
Memory cap / demand & $\mathrm{RAM}_i$, $\mathrm{RAM}^f$ & node memory; per-replica memory of $f$ [MB].\\
Service demand & $D_i^f$ & average execution time of $f$ at $i$ [s].\\
Max utilization & $U^f_{\max}$ & stability threshold, $U^f_{\max}\in[0,1)$.\\
Local / offload / Cloud profit & $\alpha_i^f$, $\beta_{ij}^f$, $\gamma_i^f$ & profit of local exec.; profit of offloading $i\!\to\! j$; Cloud-offloading penalty (per req/s).\\
Communication latency & $L_{ij}$ & horizontal-offloading latency between $i$ and $j$.\\
Weights & $w_{\mathrm{lat}}$, $w_{\mathrm{fair}}$ & importance of latency and of fairness.\\
Max iterations & $H_{\max}$ & auction iterations per control period.\\
Price params \emph{(MADeA only)} & $\varepsilon$, $\zeta$ & min.\ bid increment; idle-seller price discount $\zeta\in(0,1)$.\\
\midrule
\multicolumn{3}{@{}l}{\textit{Variables and state at iteration $h$}}\\
Replicas & $r_i^f$ & number of $f$ replicas on $i$, $r_i^f\in\mathbb{N}_0$.\\
Local / forwarded / Cloud load & $x_i^f$, $y_{ij}^f$, $z_i^f$ & locally served; forwarded $i\!\to\! j$; Cloud-offloaded [req/s].\\
Offloaded demand & $\omega_i^f$ & horizontal demand still to place [req/s].\\
Effective load & $\ell_i^f$ & load to be served at $i$ for $f$ [req/s].\\
Capacity / residual / slack & $\mathrm{Cap}_i^f$, $C_i^f$, $\rho_i$ & exec.\ capacity (Eq.~2); residual exec.\ capacity (Eq.~12); memory slack (Eq.~13).\\
Congestion ratio & $\kappa_i^f$ & utilization ratio for $f$ on $i$, $\kappa_i^f\in[0,1]$.\\
Fairness term & $\phi_i^f$ & penalty on nodes that offload excessively.\\
Execution price \emph{(MADeA only)} & $p_i^f$ & shadow cost of residual capacity; \textbf{removed} in \textsc{FaaS-MADiG}.\\
Coordination score & $s_{ij}^f$ & price-free preference of buyer $i$ for seller $j$ (Eq.~\eqref{eq:madig-score}).\\
\bottomrule
\end{tabular}
\end{table}
% ===================== END self-contained notation recap =====================

\subsection{Price-free coordination score}
\label{sec:madig-score}

At iteration $h$, an overloaded buyer $i$ (with $\omega_i^f(h)>0$) ranks each
neighbour $j$ advertising residual capacity $C_j^f(h)\ge 1$ by the
\emph{same} localized utility of Eq.~(7) with the execution price removed:
\begin{equation}
  s_{ij}^f(h)\;=\;u_{ij}^f(h)+p_j^f(h)
  \;=\;\beta_{ij}^f\;-\;w_{\mathrm{lat}}\,L_{ij}\;-\;w_{\mathrm{fair}}\,\phi_i^f(h),
  \label{eq:madig-score}
\end{equation}
and retains the same convenience threshold used by \textsc{FaaS-MADeA},
\begin{equation}
  j\in A_i^f(h)\iff s_{ij}^f(h)>-\,\gamma_i^f,
  \label{eq:madig-threshold}
\end{equation}
where $\beta_{ij}^f$ is the reward of serving $i$'s load for $f$ at $j$,
$L_{ij}$ the one-hop network latency, $\phi_i^f(h)$ the fairness penalty, and
$\gamma_i^f$ the vertical (Cloud) offloading cost that sets the
break-even point against Cloud forwarding. Dropping the $-p_j^f(h)$ term
removes the only quantity the auction feeds back \emph{across} iterations;
every other component of the buyer's preference, and the seller it prefers, are
left untouched.

\subsection{Greedy diffusion}
\label{sec:madig-greedy}

Two deterministic rules replace the auction's bid/clear cycle. On the buyer
side (Algorithm~\ref{alg:madig-buyer}), node $i$ sorts its candidate sellers by
descending score~\eqref{eq:madig-score} and greedily claims residual capacity
until its demand $\omega_i^f(h)$ is met or the neighbourhood is saturated; any
shortfall triggers \emph{price-free} replica-expansion requests to neighbours
with memory slack $\rho_j(h)>0$, exactly as in \textsc{FaaS-MADeA}. On the
seller side (Algorithm~\ref{alg:madig-seller}), node $j$ serves incoming
requests in descending-score order---ties broken by the buyer index $i$ for
reproducibility---up to its residual capacity $C_j^f(h)$. If capacity is
exhausted while expansion requests remain and memory permits, the seller
\emph{tentatively} starts replicas through the unchanged mechanism of
Algorithm~4, which is gated by the utilization test $\kappa_j^f\le 1$ (i.e.
$u\le U^f_{\max}$) and therefore needs no price.

\begin{algorithm}
\caption{\textsc{FaaS-MADiG} --- buyer greedy assignment (replaces Alg.~2)}
\label{alg:madig-buyer}
\begin{algorithmic}[1]
\Require residual demand $\omega_i^f(h)$, neighbour residuals $C_j^f(h)$ and
         memory slacks $\rho_j(h)$ on the blackboard
\For{each function $f$ with $\omega_i^f(h)>0$}
  \State $A_i^f(h)\gets\{\,j\in N_i: C_j^f(h)\ge 1 \text{ and } s_{ij}^f(h)>-\gamma_i^f\,\}$
  \State sort $A_i^f(h)$ by \textbf{descending} $s_{ij}^f(h)$
  \State $o\gets 0$ \Comment{load already placed}
  \For{$j\in A_i^f(h)$ \textbf{while} $o<\omega_i^f(h)$}
     \State $d\gets\min\!\big(C_j^f(h),\,\omega_i^f(h)-o\big)$ \Comment{unit bids if enabled}
     \State emit assignment request $(i\!\to\! j,f,d)$;\quad $o\gets o+d$
  \EndFor
  \If{$o<\omega_i^f(h)$} \Comment{capacity insufficient}
     \For{$j\in N_i$ with $\rho_j(h)>0$ and $j$ not already a capacity seller}
        \State emit \emph{price-free} replica-expansion request $(i\!\to\! j,f)$
     \EndFor
  \EndIf
\EndFor
\end{algorithmic}
\end{algorithm}

\begin{algorithm}
\caption{\textsc{FaaS-MADiG} --- seller greedy fill (replaces Alg.~3)}
\label{alg:madig-seller}
\begin{algorithmic}[1]
\Require requests received by $j$, residual capacity $C_j^f(h)$, memory slack $\rho_j(h)$
\For{each function $f$ with $C_j^f(h)>0$ \textbf{or} pending expansion requests}
  \State sort received requests by $\big(\text{descending } s_{ij}^f(h),\ \text{ascending } i\big)$
  \State $c\gets C_j^f(h)$
  \While{requests remain \textbf{and} $c>0$}
     \State $q\gets\min(c,\,d_{\text{req}})$;\quad $y_{ij}^f\mathrel{+}= q$;\quad $c\gets c-q$
  \EndWhile
  \If{$c=0$ \textbf{and} expansion requests remain \textbf{and} $\rho_j(h)>0$}
     \State tentatively start replicas as in Alg.~4 \Comment{utilization test $\kappa_j^f\le1$; no price}
  \EndIf
\EndFor
\State \textbf{no} price update, \textbf{no} bid-price tracking, \textbf{no} reassignment of previously served load
\end{algorithmic}
\end{algorithm}

\subsection{What is removed and what is kept}
\label{sec:madig-ablation}

\noindent\textbf{Removed} (the price machinery of Eqs.~(7)--(8) and
Algorithms~2--3): the price term $-p_j^f(h)$ in the utility; the per-unit
bidding price $b_{ij}^f$ and minimum increment $\varepsilon$; the
congestion-based price update of Eq.~(8) and the idle-seller discount $\zeta$;
and the price-gated reassignment of load already granted to another buyer.

\noindent\textbf{Kept} (identical to \textsc{FaaS-MADeA}): the local
problem~(P2); the advertising of residual execution capacity $C_i^f(h)$ and
memory slack $\rho_i(h)$ (Eqs.~(12)--(13)); the convenience threshold
$\gamma_i^f$ and the fairness penalty $\phi_i^f$; the price-free replica
expansion of Algorithm~4; and the restricted-problem re-solve after each
coordination round.

\subsection{Rationale and properties}
\label{sec:madig-rationale}

\paragraph{Clean ablation.} Because the buyer already ranks sellers greedily by
utility in \textsc{FaaS-MADeA}, the \emph{only} behavioural difference of
\textsc{FaaS-MADiG} is the absence of the price signal---the cross-iteration
congestion feedback and the seller-side price arbitration. Any performance gap
between the two is therefore attributable to that signal alone, which makes
\textsc{FaaS-MADiG} a controlled baseline rather than an unrelated heuristic.

\paragraph{Determinism and termination.} With prices removed, the score
\eqref{eq:madig-score} is constant within a control period $t$, the
tie-break on the buyer index makes the seller fill deterministic, and residual
capacity is consumed monotonically within each iteration. The stopping criteria
of \textsc{FaaS-MADeA} carry over unchanged (no residual capacity, all load
assigned, $H_{\max}$ reached, or no convenient seller). Removing the
price-gated reassignment also removes the main source of inter-iteration
oscillation, so a price-free run need not damp prices to stabilize.

\paragraph{Fair runtime comparison.} Heuristic time is amortized per agent
(the per-iteration cost is divided by the number of participating
buyers/sellers, as in \textsc{FaaS-MADeA}), and only the MILP solves are charged
in full. Wall-clock differences between the two methods thus reflect the cost of
the coordination mechanism, not of the implementation.

\paragraph{Communication.} \textsc{FaaS-MADiG} uses the same one-hop blackboard
as \textsc{FaaS-MADeA} (advertising $C_i^f(h)$ and $\rho_i(h)$), but exchanges
\emph{no} prices and runs \emph{no} price-update round, so its per-iteration
message footprint is strictly smaller.

\subsection{Positioning with respect to the literature}
\label{sec:madig-related}

We make explicit that \textsc{FaaS-MADiG} is, by design, an instance of a
well-established family of distributed load-balancing rules, used here as a
\emph{controlled} baseline rather than as a novel algorithm. Stating its lineage
clarifies that the contribution it supports is methodological---isolating the
effect of congestion pricing---and not a new heuristic.

\paragraph{Similarities.} The buyer-side push of residual demand to one-hop
neighbours ranked by a local gradient is exactly the \emph{diffusion} paradigm of
Cybenko~\citep{cybenko1989}, in which an overloaded node iteratively sends excess
load to its less-loaded neighbours; the sender-initiated form we use is a
single-round simplification of the sender-initiated diffusion (SID) strategy
catalogued by Willebeek-LeMair and Reeves~\citep{willebeek1993}. Greedy and
heuristic offloading is likewise common in Edge/FaaS load management---for
instance greedy, locality-driven load balancing inside FaaS
dispatchers~\citep{leechoi2021} and decentralized multi-agent edge load
balancing~\citep{nezami2021}---although those works differ in mechanism
(dispatcher-level cache-locality scheduling, and collective-learning plan
selection with randomly sampled candidate hosts, respectively) rather than
instantiating the same one-hop residual-offload rule. Finally, the relation to
\textsc{FaaS-MADeA} is precise: \textsc{FaaS-MADiG} is the \emph{price-frozen}
special case of a Bertsekas-style auction~\citep{bertsekas1988}. Removing the
cross-iteration price update---together with the $\varepsilon$ bid increment and
the $\varepsilon$-complementary-slackness mechanism it serves---collapses the
bidding-and-clearing cycle of
Algorithms~\ref{alg:madig-buyer}--\ref{alg:madig-seller} to a single greedy pass
in which each buyer claims its best available residual capacity and \emph{no}
previously served load is re-awarded. In this sense \textsc{FaaS-MADiG}
\emph{is} \textsc{FaaS-MADeA} with the market layer removed.

\paragraph{Differences.} Three elements separate \textsc{FaaS-MADiG} from the
textbook rules above. \emph{(i) Determinism over the full neighbourhood.} Unlike
randomized power-of-$d$-choices balancing~\citep{mitzenmacher2001}, which probes
a random sample of $d$ candidates, \textsc{FaaS-MADiG} scans the entire one-hop
candidate set $A_i^f(h)$ and breaks ties deterministically, yielding reproducible
runs without seed averaging. \emph{(ii) Joint balancing and provisioning.}
Classical diffusion migrates load over a \emph{fixed} resource graph;
\textsc{FaaS-MADiG} additionally triggers memory-slack-based replica expansion
(Algorithm~4), coupling load balancing with capacity provisioning---a feature
inherited from the \textsc{FaaS-MADeA} framework rather than from diffusion.
\emph{(iii) Objective-aligned, optimization-coupled ranking.} The score
\eqref{eq:madig-score} ranks neighbours by the model's economic weights
($\beta_{ij}^f$, $\gamma_i^f$, latency $L_{ij}$, fairness $\phi_i^f$) rather than
by raw queue length, and every offloading round is re-validated by the local
problem~(P2) and the restricted re-solve, so the diffusion is steered by the
social-welfare objective instead of by local load alone.

\paragraph{Takeaway.} As an isolated coordination rule, \textsc{FaaS-MADiG}
instantiates the diffusion / greedy one-hop-offloading paradigm and is not
claimed as novel; its value is as a price-free ablation that attributes any
performance gap with respect to \textsc{FaaS-MADeA} to congestion pricing alone.
